\section*{39. 超复系统及其与交换代数和数论的关系}

\begin{center}
\emph{Verhandl. Intern. Math.-Kongreß Zürich 1 (1932), S. 189--194}
\end{center}

\begin{center}
\textsc{由 Emmy Noether 撰，哥廷根}
\end{center}

\paragraph{1.}
超复系统，也就是代数的理论，在近几年有了强劲发展；但是直到最近，这一理论对于交换性问题的意义才变得清楚。今天我要报告非交换理论对于交换理论的这种意义；具体说，我将沿着两个可追溯到 Gauss 的经典问题来说明，即主属定理以及与它密切相关的范数定理。这些问题在历史进程中不断改变表述形式：在 Gauss 那里，它们作为二次型理论的收尾出现；随后，它们在通过类域论刻画相对循环数域和阿贝尔数域时起到本质作用；最后，它们可以表述为关于自同构以及关于代数分裂的定理，而这一最后表述同时把这些定理推广到任意相对伽罗瓦数域。

通过这幅稍后还要展开的草图，我同时想说明把非交换理论应用于交换理论的\emph{原则}：\emph{借助代数理论，人们力图为关于二次型或循环域的已知事实取得不变而简单的表述，也就是说，取得只依赖于代数结构性质的表述。一旦这些不变表述已经得到证明，而且上述例子正是如此，那么这些事实向任意伽罗瓦域的转移也就自动获得。}

\paragraph{2.}
在展开细节之前，我还要先概述各种方法和进一步结果。首先必须指出，主要困难在于为一般伽罗瓦域取得相应表述；若没有超复方法，原本甚至没有着手点。在上面列举的例子中，向非交换理论的根本过渡，是通过\emph{同时考察域与群}而取得的，具体说，是借助“交叉积”及其乘法常数，也就是“因子系”（见 3）。这样得到一个关于基域的单纯正规代数，而每一个这样的代数本质上也都能这样生成。这样的交叉积最早由 Dickson 考察，\footnote{参见他的著作 \emph{Algebren und ihre Zahlentheorie}, Zürich 1927, § 34。}而因子系理论则由 Speiser、Schur 和 R. Brauer\footnote{例如参见 R. Brauer, \emph{Untersuchungen über die arithmetischen Eigenschaften von Gruppen linearer Substitutionen}, 以及该文注 2 所列文献，Math. Z. 28 (1928)。}从完全不同的出发点发展出来，即从绝对不可约表示的问题出发。只有通过这两种理论的融合，才可能得到一个足够简单且影响范围足够广的结构，从而能够用它处理交换性问题。\footnote{这一结构最初由我在 1929/30 年冬季学期的一门课程中发展出来，后来收入 H. Hasse, \emph{Theory of cyclic algebras}, Trans. Amer. Math. Soc. 34 (1932), 第 2 章。关于本报告涉及的整个领域，M. Deuring 关于超复数及其数论应用的一篇报告将提供综述，该报告拟收入 \emph{Ergebnisse der Mathematik} 丛书。}

与此同时，对一些已知事实也得到新的、透明的证明：这里我想指出 H. Hasse 给出的循环域互反律的超复证明，这一证明即将发表于 Math. Ann.，\footnote{H. Hasse, \emph{Die Struktur der R. Brauerschen Algebrenklassengruppe über einem algebraischen Zahlkörper (insbesondere Normenrestsymbol und Reziprozitätsgesetz)}. Math. Ann. 107 (1932/33)。}它依靠的是基于交叉积理论的范数剩余符号的不变表述。还要指出 C. Chevalley 最近在同一基础上给出的局部类域论的超复奠基；不过为此还必须发展若干关于因子系的新代数定理。\footnote{C. Chevalley, \emph{Sur la théorie du symbole de restes normiques}; 将发表于 J. f. Math. 169。}同时我还必须作一个限制性说明：只靠交叉积方法，看来并不能得到伽罗瓦数域的完整理论。这一点来自 Artin 尚未发表的新结果；这些结果按照上述原则接续 Hasse 的证明，但只能给出计数相等，而不是完整的同构定理。

能够给出完整同构，准确说是算子同构的方法，已经存在于代数方面。这里涉及的是 A. Speiser\footnote{A. Speiser, \emph{Gruppendeterminante und Körperdiskriminante}. Math. Ann. 77 (1916)。}若干思路的延续，即把伽罗瓦域理解为“伽罗瓦模”：也就是说，把它看成基域上的一个模，其中伽罗瓦群的代换作为算子出现。于是域与群环（群代数）之间存在算子同构，其意义是：元素之间有一一对应，基域上的线性形式相互对应，并且域中伽罗瓦群的代换对应于群环中的乘法。我提出的这一定理\footnote{E. Noether, \emph{Normalbasis bei Körpern ohne höhere Verzweigung}, Satz 3, J. f. Math. 167 (1932)（证明中含有一处缺口）。}已经由 M. Deuring 证明，\footnote{M. Deuring, \emph{Galoissche Theorie und Darstellungstheorie}. Math. Ann. 107 (1932)。}他在此基础上给出了伽罗瓦理论的一个证明，其中算子同构实现了群与域的对应。进一步的结构定理，同样由 Deuring 给出，与 Artin \(L\)-级数的形式事实相平行，并给出通向 Artin 导子的结构性途径。这些 Artin \(L\)-级数和导子\footnote{E. Artin, \emph{Über eine neue Art von \(L\)-Reihen}, Math. Sem. Hamburg 3 (1924). \emph{Zur Theorie der \(L\)-Reihen mit allgemeinen Gruppencharakteren}, 同刊, 8 (1931). \emph{Die gruppentheoretische Struktur der Diskriminanten algebraischer Zahlkörper}, J. f. M. 164 (1931).}由一般群特征标构成，和 Speiser\textsuperscript{6)} 的工作一起，是数论与表示论之间的第一次联系，是越过阿贝尔域范围的第一次推进。它们给整个发展以强烈推动；特别是伽罗瓦模理论正是以此为取向发展起来的。

\paragraph{3.}
现在我要具体追踪置于开头的两个问题，即范数定理和主属定理。首先是\emph{交叉积的定义}：设 \(K/k\) 是一个 \(n\) 次伽罗瓦域，\(\Gg\) 为其群。所谓交叉积，是指把 \(K\) 和 \(\Gg\) 同时嵌入一个代数 \(A\)，使得 \(K\) 的自同构成为内自同构。令 \(u_{S_1},\ldots,u_{S_n}\) 表示与这 \(n\) 个群元素相应的符号；于是先把 \(A\) 设为关于 \(K\) 的秩 \(n\) 线性形式模：
\[
\tag{1}
A=u_{S_1}K+\cdots+u_{S_n}K .
\]
也就是说，\(A\) 由所有线性形式 \(u_{S_1}a_1+\cdots+u_{S_n}a_n\) 组成，其中 \(a_i\) 在 \(K\) 中任意取值。通过要求由 \(u_S\)，更一般地由 \(u_SK^\ast\)\footnote{\(K^\ast\) 由 \(K\) 去掉零元而得；这一记号通用。}生成内自同构，\(A\) 便成为一个环，也就是一个关于 \(k\) 的秩 \(n^2\) 代数。这个要求表示为
\[
\tag{2}
u_S^{-1}zu_S=z^S,\qquad \text{或}\qquad zu_S=u_Sz^S
\]
对每个 \(z\in K\) 成立，\footnote{\(z^S\) 照常表示由代换 \(S\) 作用于 \(z\) 后所得的元素。}并且乘法满足
\[
\tag{3}
u_Su_T=u_{ST}a_{S,T}\quad \text{其中 }a_{S,T}\text{ 在 }K^\ast\text{ 中},
\]
\[
\tag{4}
\begin{gathered}
a_{ST,R}\,a_{S,T}^{\,R}=a_{S,TR}\,a_{T,R}\\
\text{（由结合律 \([u_Su_T]u_R=u_S[u_Tu_R]\) 得到）}.
\end{gathered}
\]
\(A\) 称为 \(K\) 与 \(\Gg\) 在因子系 \(a_{S,T}\) 下的交叉积；可以证明，\(A\) 成为 \(k\) 上的一个单纯正规代数，也就是某个全矩阵环 \(D_r\)，其系数取自在相应除代数 \(D\) 中，并且 \(K\) 成为极大交换子域，因而成为分裂域（即把系数域 \(k\) 扩张到一个与 \(K\) 同构的域后，所得代数分裂为以中心为系数的全矩阵环）。反过来，对预先给定的除代数 \(D\)，总存在矩阵环 \(D_r\)，可按上述方式生成作为交叉积。

如果从 \(u_S\) 过渡到 \(v_S=u_Sc_S\)，其中 \(c_S\) 在 \(K^\ast\) 中，这仍生成同一个自同构，但会产生“相伴”的因子系：
\[
\tag{5}
\bar a_{S,T}=a_{S,T}\frac{c_S^{\,T}c_T}{c_{ST}} .
\]
相伴的因子系合并为一个类 \((a)\)；同样，把所有与 \(A\) 相似的代数，也就是所有 \(r=1,2,\ldots\) 的 \(D_r\)，合并为一个类 \(\mathfrak A\)。\emph{类 \(\mathfrak A\) 与类 \((a)\) 一一对应：具有固定分裂域 \(K\) 的这些类，在直接乘法下构成一个阿贝尔群，它同因子系类的逐项乘积群同构。单位元是分裂代数类，或者说是所有变换量所构成的系统}
\[
        \frac{c_S^{\,T}c_T}{c_{ST}} .
\]
这里涉及的是 R. Brauer 所指出的代数类群。

\paragraph{4.}
现在我通过\emph{特殊化}到\emph{循环分裂域}来说明同\emph{范数概念}的联系，并由此按照开头说明的原则得到推广\emph{范数定理}的表述。若 \(Z\) 是循环域，\(S\) 是其群的一个生成代换，那么可以让 \(S\) 的幂对应于 \(u\) 的幂，即
\[
\begin{aligned}
(1')\qquad &A=Z+uZ+\cdots+u^{n-1}Z\\
(2')\qquad &zu=uz^S\\
(3')\qquad &u^n=a\\
(4')\qquad &a\text{ 位于基域 }k^\ast\text{ 中}\\
(5')\qquad &\bar a=a\cdot \N(c),\quad \text{当置 }v=uc\text{ 时}.
\end{aligned}
\]
因此这里每一个因子系只由一个位于基域中的元素 \(a\) 组成，记作 \(A=(a,Z)\)；因子系的单位类由来自 \(Z^\ast\) 的范数给出，而代数类群同商群
\[
        k^\ast/\N(Z^\ast)
\]
同构。于是循环代数 \((a,Z)\) 当且仅当 \(a\) 是 \(Z\) 中某个元素的范数时分裂。

范数与分裂之间的这一联系给出了“范数定理”的表述，即\emph{关于分裂代数的定理：若一个代数在每个素位处分裂，则它本身分裂。}这里的“素位”按数论中的通常方式定义：把基域 \(k\) 替换为它的 \(p\)-进扩张 \(k_{\mathfrak p}\)，其中 \(\mathfrak p\) 是 \(k\) 的一个素理想；对有限多个无限素位，则相应地把 \(k\) 及其共轭域用实数域来扩张。

事实上，循环域的范数定理包含在这一定理之中。因为对循环代数 \((a,Z)\) 而言，根据上面的说明，这一定理等价于如下断言：若 \(a\) 在每个有限或无限素位处都是 \(p\)-进范数，则 \(a\) 是 \(Z\) 中某个数的范数；或者不经过 \(p\)-进过渡地说：\emph{若 \(a\) 按 \(k\) 的每个素理想 \(\mathfrak p\) 都是范数剩余（并满足某些符号条件），则 \(a\) 是 \(Z\) 中某个数的范数。}后一种表述正是类域论中利用熟知的解析工具所证明的范数定理。而关于分裂代数的一般定理，则可由这一循环特殊情形通过纯代数-算术的考察得到。\footnote{R. Brauer, H. Hasse, E. Noether, \emph{Beweis eines Hauptsatzes in der Theorie der Algebren}. J. f. M. 167 (1932)。}Hasse 已由此得出第一个重要推论：\emph{代数数域上的每个单纯正规代数都是循环的。}正是在寻找这个长期被猜测事实的证明时，才产生了上述一般表述。

\paragraph{5.}
从关于分裂代数的定理还可推出第二个后果，即开头提出的主属定理，仍然只需纯代数-算术的推理。\footnote{该证明拟发表于 Math. Ann.。}它的不变表述基于这样一个事实：定义交叉积的关系 (2) 到 (5) 完全是乘法性的；因此，如果把 \(K^\ast\) 替换为一个阿贝尔群 \(\mathfrak J\)，只要该群的自同构群含有一个同 \(\Gg\) 同构的子群，这些关系仍有意义。于是取代 (1) 的，是群论意义下“用 \(\mathfrak J\) 扩张 \(\Gg\)”。若把 \(\mathfrak J\) 取为 \(K\) 中所有理想构成的群，因子系便成为理想系统；\(\mathfrak J\) 中的一个类划分诱导出因子系的一个类划分，而且一般会得到比原来的理想类划分更细的划分。因为要求 \(u_S\) 的乘法对绝对理想类是单值的，只说变换量 \(c_S^{\,T}c_T/c_{ST}\)，也就是元素因子系的单位类，落在理想因子系的单位类中。事实上，已知特殊情形表明，稍微不那么细的划分已经足够。我作如下定义：\emph{把那些在 \(K\) 的所有有限和无限分歧素位处生成分裂代数的元素 \(a_{S,T}\)，计入因子系的单位类。}

这样产生的 \(\Gg\) 的扩张记作 \(\Gg^\ast\)；\((c)\) 表示 \(c\) 的绝对理想类。于是：

\emph{主属定理的不变表述为：在这样定义的类划分下，若通过代换 \(v_S=u_S(c_S)\) 产生 \(\Gg^\ast\) 的一个自同构，其中所有这些 \((c_S)\) 构成主属，那么这个自同构是内的，并且由一个理想类 \((\mathfrak b)\) 生成。}

已知特殊情形可由一个等价而稍显明白的表述推出：\emph{若由理想类 \((c_S)\) 形成的变换量 \((c_S^{\,T})(c_T)/(c_{ST})\) 属于因子系的单位理想类，则存在一个理想类 \((\mathfrak b)\)，使得 \((c_S)\) 成为符号意义下的 \((1-S)\) 次幂：对 \(\Gg\) 中所有 \(S\)，都有 \((c_S)=(\mathfrak b)/(\mathfrak b^{S})\)。}因为前提恰好表达了自同构性质；而这个自同构是内的这一事实，则由
\[
        v_S=(\mathfrak b)^{-1}u_S(\mathfrak b)
        =u_S(\mathfrak b)^{1-S}=u_S(c_S)
\]
来表达。

特殊化到循环情形就得到（注意归一化）：若 \(N([c])\) 位于因子系的单位理想类中，则 \((c)\) 成为符号意义下的 \((1-S)\) 次幂：
\[
        (c)=(\mathfrak b)^{1-S}.
\]

\paragraph{6.}
为了由此转到循环域和二次型的熟知表述，我首先指出，定理虽然是针对完整理想类说出的，但如果像通常那样只限于同分歧素位互素的理想，它的内容仍保持不变。

这样，这里定义的二次域主属就转化为 Gauss 的主属；因为理想类对应于二次型，而类的范数对应于可由这些型表示的数。单位类由在 \(K\) 的分歧素位处生成分裂代数的那些数给出，这就是说，这些可表示的数在分歧素位处成为二次剩余；相应的型因而具有主型的全特征，所以构成 Gauss 的主属。符号意义下的 \((1-S)\) 次幂转化为平方化，这是熟知的。

对于循环域，表述转化为如下形式：主属由所有理想类组成，其范数在有限和无限分歧素位处都成为范数剩余。但这又等价于熟知的“按导子的范数剩余”；通常的定理便这样作为特殊化得到。

此外，即使在任意伽罗瓦域的一般情形中，也可以引入一个只由分歧素位组成的导子，使得在对因子系作归一化时，对每个这样的素位，射线只含有单位类中的元素。

这里产生的问题是：这同概述 2 中提到的 Artin 导子有什么联系？那些导子毕竟也是由同一些素理想组成的；因此也产生了它同伽罗瓦模理论，即第二种超复方法之间有什么联系的问题。这两种方法究竟能推进到何种程度，还有待将来说明。

% END live source-fidelity unit: translations/non_slavic/simplified_chinese/paper39/source_fidelity/v001/Noether_Paper39_SourceFidelity_SimplifiedChinese_v001.tex
\clearpage
% BEGIN live source-fidelity unit: translations/non_slavic/simplified_chinese/paper40/source_fidelity/v001/Noether_Paper40_SourceFidelity_SimplifiedChinese_v001.tex
\setcounter{footnote}{0}

% Source-fidelity v001: Paper40 complete Simplified Chinese translation unit.
% Authority: repaired R124+P40 Math. Z. 37 witness.
% Source witness: sources/paper40/source_fidelity/Noether_Paper40_ORIGINAL_MathZ37_R124plusP40_repaired_witness_v001.tex.
% Boundary: older German/English local slices are retained as controls only; translation follows the repaired witness through the received-date line.
